\section*{CHƯƠNG 3. CƠ SỞ LÝ THUYẾT VÀ PHƯƠNG PHÁP}
\setcounter{section}{3}
\setcounter{subsection}{0}
\addcontentsline{toc}{section}{Chương 3. Cơ sở lý thuyết và Phương pháp}

Chương này trình bày nền tảng toán học và giải thuật của bốn nhóm phương pháp tương ứng với các notebook thực nghiệm \code{1\_Bayes}, \code{2\_Logistic}, \code{3\_MLP\_CNN} và \code{4\_Ensemble}.

\subsection{Mô hình xác suất Naïve Bayes đa thức (\code{MultinomialNB-NP})}
\subsubsection{Cơ sở toán học}
Mô hình Naïve Bayes dựa trên định lý Bayes để tính xác suất hậu nghiệm của lớp $y = c$ khi quan sát vector điểm ảnh $\mathbf{x} = [x_1, \dots, x_D]^T \in \mathbb{R}^{784}$:
\begin{equation}
    P(y = c \mid \mathbf{x}) \propto P(y = c) P(\mathbf{x} \mid y = c) = P(y = c) \prod_{i=1}^D P(x_i \mid y = c)
\end{equation}

Trong biến thể Naïve Bayes đa thức, mỗi giá trị $x_i \in [0, 1]$ được xem như cường độ xuất hiện của đặc trưng $i$. Xác suất có điều kiện được mô hình hóa qua tham số $\theta_{ci}$ kết hợp kỹ thuật làm trơn Laplace ($\alpha > 0$):
\begin{equation}
    \theta_{ci} = \frac{N_{ci} + \alpha}{N_c + \alpha D}, \qquad N_{ci} = \sum_{j: y^{(j)}=c} x_{j, i}, \quad N_c = \sum_{i=1}^D N_{ci}
\end{equation}

Quy tắc phân loại cực đại hậu nghiệm (MAP) trong không gian logarit có dạng:
\begin{equation}
    \hat{y} = \arg\max_{c \in \{0, \dots, 9\}} \left[ \ln P(y = c) + \sum_{i=1}^D x_i \ln \theta_{ci} \right]
\end{equation}

\begin{lstlisting}[style=pseudocode, caption={Thuật toán Multinomial Naïve Bayes (Laplace Smoothing)}]
ALGORITHM: Multinomial Naive Bayes (Fit & Predict)
INPUT: X in R^(N x D), y in {0..K-1}^N, he so alpha > 0, X_test in R^(M x D)
OUTPUT: Nhan du doan y_pred cho X_test
1: FOR c = 0 TO K - 1 DO:
2:   prior[c] = count(y == c) / N
3:   N_ci = sum(X[y == c], axis=0)                // Tong cuong do pixel lop c
4:   Theta[c] = (N_ci + alpha) / (sum(N_ci) + alpha * D)
5: ENDFOR
6: Scores = ln(prior) + X_test @ ln(Theta.T)      // Tich vo huong log-likelihood
7: RETURN argmax(Scores, axis=1)
\end{lstlisting}

\subsection{Hồi quy Logistic đa lớp (\code{LogisticRegression-NP})}
\subsubsection{Cơ sở toán học}
Mô hình Hồi quy Logistic đa lớp được thiết lập theo chiến lược Một-đấu-Tất cả (\textbf{One-Vs-Rest} - OvR). Với bài toán $K = 10$ lớp, mô hình huấn luyện độc lập $10$ bộ phân loại nhị phân. Hàm Sigmoid ánh xạ tổ hợp tuyến tính về miền xác suất:
\begin{equation}
    \hat{p}_k(\mathbf{x}) = \sigma(\mathbf{w}_k^T \mathbf{x} + b_k) = \frac{1}{1 + e^{-(\mathbf{w}_k^T \mathbf{x} + b_k)}}
\end{equation}

Hàm mất mát Cross-Entropy nhị phân kết hợp điều chuẩn $L_2$ (không phạt bias $b_k$):
\begin{equation}
    \mathcal{J}(\mathbf{w}_k, b_k) = -\frac{1}{N} \sum_{j=1}^N \left[ y_k^{(j)} \ln \hat{p}_k^{(j)} + (1 - y_k^{(j)}) \ln (1 - \hat{p}_k^{(j)}) \right] + \frac{\lambda}{2} \|\mathbf{w}_k\|_2^2
\end{equation}

Gradient giải tích và quy tắc cập nhật Gradient Descent:
\begin{equation}
    \nabla_{\mathbf{w}_k} \mathcal{J} = \frac{1}{N} \mathbf{X}^T (\hat{\mathbf{p}}_k - \mathbf{y}_k) + \lambda \mathbf{w}_k, \quad \nabla_{b_k} \mathcal{J} = \frac{1}{N} \sum_{j=1}^N (\hat{p}_k^{(j)} - y_k^{(j)})
\end{equation}
\begin{equation}
    \mathbf{w}_k \leftarrow \mathbf{w}_k - \eta \nabla_{\mathbf{w}_k} \mathcal{J}, \qquad b_k \leftarrow b_k - \eta \nabla_{b_k} \mathcal{J}
\end{equation}

\begin{lstlisting}[style=pseudocode, caption={Thuật toán One-Vs-Rest Logistic Regression}]
ALGORITHM: One-Vs-Rest Logistic Regression with L2 Regularization
INPUT: X in R^(N x D), y in {0..K-1}^N, toc do hoc eta, epoch E, lambda, X_test
OUTPUT: Nhan du doan y_pred cho X_test
1: Khoi tao W = zeros(D, K), b = zeros(K)
2: FOR k = 0 TO K - 1 DO:
3:   y_k = (y == k).astype(float)
4:   FOR epoch = 1 TO E DO:
5:     p = 1.0 / (1.0 + exp(-(X @ W[:, k] + b[k])))
6:     grad_w = (X.T @ (p - y_k)) / N + lambda * W[:, k]
7:     grad_b = sum(p - y_k) / N
8:     W[:, k] -= eta * grad_w;  b[k] -= eta * grad_b
9:   ENDFOR
10: ENDFOR
11: RETURN argmax(X_test @ W + b, axis=1)
\end{lstlisting}

\subsection{Mô hình Học sâu: Mạng nơ-ron đa tầng (\code{MLP}) và Mạng tích chập (\code{CNN})}
\subsubsection{Mạng nơ-ron truyền thẳng đa tầng (\code{MLP})}
Kiến trúc MLP gồm $2$ tầng ẩn: $784 \rightarrow 256 \rightarrow 128 \rightarrow 10$. Hàm kích hoạt ReLU $\mathbf{a}^{(l)} = \max(0, \mathbf{z}^{(l)})$ đi kèm phép khởi tạo trọng số He Normal: $W^{(l)} \sim \mathcal{N}(0, 2/k_{in})$. Tầng đầu ra sử dụng hàm Softmax và hàm mất mát Cross-Entropy:
\begin{equation}
    \hat{y}_c = \frac{e^{z_c^{(L)}}}{\sum_{j=0}^{K-1} e^{z_j^{(L)}}}, \qquad \mathcal{L} = -\frac{1}{B} \sum_{j=1}^B \sum_{c=0}^{K-1} y_{jc} \ln \hat{y}_{jc}
\end{equation}

Gradient sai số tầng ra $\boldsymbol{\delta}^{(L)} = \frac{1}{B} (\hat{\mathbf{Y}} - \mathbf{Y})$. Gradient truyền ngược qua các tầng ẩn:
\begin{equation}
    \nabla_{W^{(l)}} \mathcal{L} = (\boldsymbol{\delta}^{(l)})^T \mathbf{a}^{(l-1)} + \lambda_{wd} W^{(l)}, \quad \boldsymbol{\delta}^{(l-1)} = (\boldsymbol{\delta}^{(l)} W^{(l)}) \odot \mathbb{I}(\mathbf{z}^{(l-1)} > 0)
\end{equation}

Cập nhật trọng số bằng Mini-batch SGD có quán tính Momentum ($\beta = 0.5$):
\begin{equation}
    \mathbf{v}_t = \beta \mathbf{v}_{t-1} - \eta \nabla_W \mathcal{L}, \qquad W \leftarrow W + \mathbf{v}_t
\end{equation}

\subsubsection{Mạng nơ-ron tích chập (\code{CNN})}
Kiến trúc CNN khai thác trực tiếp tương quan không gian 2D trên ma trận ảnh:
\begin{itemize}
    \item \textbf{Tầng Conv2D:} 32 bộ lọc $4 \times 4$ ($S=1, P=0$) tạo bản đồ đặc trưng $(28-4+1) \times (28-4+1) = 25 \times 25 \times 32$.
    \item \textbf{Tầng MaxPool2D:} Cửa sổ $2 \times 2$, bước trượt 2 làm giảm kích thước xuống $12 \times 12 \times 32 = 4.608$ chiều, tạo tính bất biến dịch chuyển.
    \item \textbf{Tầng Dense 128 \& Dense 10:} Chiếu phi tuyến qua 128 nơ-ron (ReLU) trước khi phân loại 10 lớp Softmax. Tối ưu bằng thuật toán Adam ($\text{lr} = 10^{-3}$).
\end{itemize}

\begin{lstlisting}[style=pseudocode, caption={Thuật toán huấn luyện MLP From-Scratch với SGD Momentum}]
ALGORITHM: MLP From-Scratch Training with Mini-batch SGD Momentum
INPUT: X_train in R^(N x 784), y_train, sizes=[784, 256, 128, 10], eta, beta, lambda_wd
1: Khoi tao W^(l) bang He Normal, b^(l) = 0, van toc v_W^(l) = 0, v_b^(l) = 0
2: FOR epoch = 1 TO epochs DO:
3:   FOR EACH mini-batch (X_b, Y_b) DO:
4:     // Lan truyen thuan: a^(0) = X_b -> z^(l) = a^(l-1) W^(l).T + b^(l) -> a^(l) = ReLU(z^(l))
5:     probs = Softmax(z^(L))
6:     delta = (probs - Y_b) / batch_size
7:     FOR l = L XUONG 1 DO:
8:       grad_W = delta.T @ a^(l-1) + lambda_wd * W^(l); grad_b = sum(delta, axis=0)
9:       IF l > 1 THEN delta = (delta @ W^(l)) * (z^(l-1) > 0) ENDIF
10:      v_W^(l) = beta * v_W^(l) - eta * grad_W;  W^(l) += v_W^(l)
11:      v_b^(l) = beta * v_b^(l) - eta * grad_b;  b^(l) += v_b^(l)
12:    ENDFOR
13:  ENDFOR
14: ENDFOR
\end{lstlisting}

\subsection{Học kết hợp Boosting (\code{AdaBoost})}
\subsubsection{Cơ sở toán học}
Thuật toán AdaBoost đa lớp (chuẩn SAMME) kết hợp tuần tự $T$ bộ phân loại yếu $h_t(\mathbf{x})$:
\begin{equation}
    H(\mathbf{x}) = \arg\max_{k \in \{0, \dots, K-1\}} \sum_{t=1}^T \alpha_t \cdot \mathbb{I}(h_t(\mathbf{x}) = k)
\end{equation}

Khởi tạo trọng số $w_1^{(j)} = \frac{1}{N}$. Tại mỗi vòng $t$, sau khi huấn luyện $h_t$, tính sai số có trọng số $\epsilon_t$ và hệ số tin cậy $\alpha_t$:
\begin{equation}
    \epsilon_t = \frac{\sum_{j=1}^N w_t^{(j)} \mathbb{I}(h_t(\mathbf{x}^{(j)}) \neq y^{(j)})}{\sum_{j=1}^N w_t^{(j)}}, \qquad \alpha_t = \ln\left(\frac{1 - \epsilon_t}{\epsilon_t}\right) + \ln(K - 1)
\end{equation}
Cập nhật trọng số mẫu: $w_{t+1}^{(j)} = w_t^{(j)} \exp(\alpha_t \mathbb{I}(h_t(\mathbf{x}^{(j)}) \neq y^{(j)}))$ và chuẩn hóa $\sum_j w_{t+1}^{(j)} = 1$.

\begin{lstlisting}[style=pseudocode, caption={Thuật toán AdaBoost đa lớp (SAMME)}]
ALGORITHM: AdaBoost Multi-Class (SAMME)
INPUT: X, y, BaseEstimator, so vong lap T, toc do hoc lr
1: w = ones(N) / N
2: FOR t = 1 TO T DO:
3:   h_t = Clone(BaseEstimator).Fit(X, y, sample_weight=w)
4:   epsilon_t = sum(w * (h_t.Predict(X) != y)) / sum(w)
5:   IF epsilon_t >= (1.0 - 1.0 / K) THEN BREAK ENDIF
6:   alpha_t = lr * (ln((1.0 - epsilon_t) / epsilon_t) + ln(K - 1.0))
7:   w = w * exp(alpha_t * (h_t.Predict(X) != y));  w /= sum(w)
8:   Luu h_t va alpha_t
9: ENDFOR
\end{lstlisting}
